\section{Motivation: From Manual Blender Tuning to Agentic Dataset Construction}
\label{sec:motivation}

The motivating engineering problem is repeated render tuning. A scene-aware synthetic dataset cannot be judged only by whether an image file exists. The object must be complete, visible, grounded, plausibly lit, correctly scaled, and consistent with the target instruction. When any of these conditions fails, the next step is rarely a single deterministic fix. The builder must inspect the render, infer the failure mode, decide whether to adjust lighting, object scale, camera distance, support-plane handling, material settings, or filtering rules, and then regenerate the sample.

This pattern reveals why ad-hoc pipelines are fragile. A script sequence can generate images, masks, meshes, renders, and train pairs, but it does not by itself encode why an artifact should be trusted. It also does not preserve enough state to explain downstream regressions. If a model trained on the resulting dataset improves viewpoint alignment but loses identity consistency, the system needs to know which dataset round, object subset, render-quality signal, or validation subgroup produced the conflict.

Agentic dataset construction is a response to this closed-loop nature. The goal is not to remove all human judgment or claim that agents can replace expert dataset design. The goal is narrower and more practical: make the observation-diagnosis-adjustment loop explicit enough that VLM/CV feedback and AI coding agents can handle part of the routine tuning, while humans retain oversight at verdict boundaries. This turns manual Blender iteration into an inspectable workflow with persistent artifacts and bounded actions.

The same motivation applies beyond Blender rendering. Object-level editing data often requires controlled changes and invariant preservation, whether the source is synthetic rendering, collected real images, or a mixture of both. The framework developed here is therefore task-neutral at the abstraction level, even though the first validation case comes from scene-aware object rotation.

The report builds on the setting of instruction-guided image editing and editing datasets, where InstructPix2Pix and MagicBrush connect natural-language instructions with paired source-target examples \cite{brooks2023instructpix2pix,zhang2023magicbrush}. It is also related to synthetic data factory work such as DatasetGAN, which uses generative models to produce labeled visual data with reduced manual annotation effort \cite{zhang2021datasetgan}. The present report differs in emphasis: it treats data construction as a stateful closed-loop workflow rather than only a dataset-generation procedure.
